% GENERATED FILE. Edit canonical lessons, then run npm run generate:books.
% canonical-source-hash: fnv1a64:690384805c58c320
% canonical-lessons: MW-C87-assi, MW-C87-nabbai, MW-C87-hajar, MW-C87-lakh, MW-C87-duniya

\chapter{Eighty, Ninety, Thousand, Hundred thousand, World}
\label{ch:mw-eighty-ninety-thousand-hundred-thousand-world}
\hlchaptermodality{87}

\begin{chapteropening}
\textbf{By the end of this chapter:} \emph{I can say eighty, ninety, a thousand, a hundred thousand and the world.}

Complete the last lesson of chapter 87: I can say eighty, ninety, a thousand, a hundred thousand and the world.
\end{chapteropening}

\section[\texorpdfstring{\mw{अस्सी} (assī)}{अस्सी (assī)}]{\mw{अस्सी} (assī) — eighty}

\label{lesson:MW-C87-assi}



\begin{quote}
Before the new one: say the Marwadi for sixty, then the Marwadi for seventy.
\end{quote}

\subsection*{You'll want to know: \mw{अस्सी}}
\textbf{\mw{अस्सी}} — \emph{assī} — \textquotedblleft{}eighty\textquotedblright{}.

\begin{cousinweb}
\textbf{\mw{अस्सी}} (\emph{assī}) is a tadbhava, a word inherited from Sanskrit and reshaped by sound change. It comes from \emph{aśīti}, \textquotedblleft{}eighty\textquotedblright{}. CDIAL gives no Marwari form in this entry (CDIAL 911); the match is through Hindi \emph{assī}.
\end{cousinweb}

\begin{grammarlens}[title={What you've built}]
One more word for this chapter.
\end{grammarlens}

\subsection*{Your turn}
\begin{itemize}
\raggedright
  \item \emph{Say it:} \emph{assī}
  \item \emph{Say it:} \emph{assī}, once more
  \item \emph{Say it:} say \emph{assī}
  \item \emph{Recall:} say the Marwadi for sixty, then the Marwadi for seventy, then say \emph{assī} again
\end{itemize}

\begin{tcolorbox}[breakable,colback=teal!4,colframe=teal!35!black,title={Before you move on}]
What does \mw{अस्सी} mean? (\textquotedblleft{}Eighty\textquotedblright{}.) Say it once more.
\end{tcolorbox}
\section[\texorpdfstring{\mw{नब्बै} (nabbai)}{नब्बै (nabbai)}]{\mw{नब्बै} (nabbai) — ninety}

\label{lesson:MW-C87-nabbai}



\begin{quote}
Before the new one: say the Marwadi for seventy, then the Marwadi for eighty.
\end{quote}

\subsection*{You'll want to know: \mw{नब्बै}}
\textbf{\mw{नब्बै}} — \emph{nabbai} — \textquotedblleft{}ninety\textquotedblright{}.

\begin{cousinweb}
\textbf{\mw{नब्बै}} (\emph{nabbai}) is a tadbhava, a word inherited from Sanskrit and reshaped by sound change. It comes from \emph{navati}, \textquotedblleft{}ninety\textquotedblright{}. CDIAL gives no Marwari form in this entry (CDIAL 6995); the match is through Hindi \emph{nabbe}.
\end{cousinweb}

\begin{grammarlens}[title={What you've built}]
One more word for this chapter.
\end{grammarlens}

\subsection*{Your turn}
\begin{itemize}
\raggedright
  \item \emph{Say it:} \emph{nabbai}
  \item \emph{Say it:} \emph{nabbai}, once more
  \item \emph{Say it:} say \emph{nabbai}
  \item \emph{Recall:} say the Marwadi for seventy, then the Marwadi for eighty, then say \emph{nabbai} again
\end{itemize}

\begin{tcolorbox}[breakable,colback=teal!4,colframe=teal!35!black,title={Before you move on}]
What does \mw{नब्बै} mean? (\textquotedblleft{}Ninety\textquotedblright{}.) Say it once more.
\end{tcolorbox}
\section[\texorpdfstring{\mw{हजार} (hajār)}{हजार (hajār)}]{\mw{हजार} (hajār) — a thousand}

\label{lesson:MW-C87-hajar}



\begin{quote}
Before the new one: say the Marwadi for eighty, then the Marwadi for ninety.
\end{quote}

\subsection*{You'll want to know: \mw{हजार}}
\textbf{\mw{हजार}} — \emph{hajār} — \textquotedblleft{}a thousand\textquotedblright{}.

\begin{cousinweb}
\textbf{\mw{हजार}} (\emph{hajār}) is a loan, usually traced to Persian \emph{hazār}, \textquotedblleft{}thousand\textquotedblright{}. Marwadi says the \emph{z} as \emph{j}.
\end{cousinweb}

\begin{grammarlens}[title={What you've built}]
One more word for this chapter.
\end{grammarlens}

\subsection*{Your turn}
\begin{itemize}
\raggedright
  \item \emph{Say it:} \emph{hajār}
  \item \emph{Say it:} \emph{hajār}, once more
  \item \emph{Say it:} say \emph{hajār}
  \item \emph{Recall:} say the Marwadi for eighty, then the Marwadi for ninety, then say \emph{hajār} again
\end{itemize}

\begin{tcolorbox}[breakable,colback=teal!4,colframe=teal!35!black,title={Before you move on}]
What does \mw{हजार} mean? (\textquotedblleft{}A thousand\textquotedblright{}.) Say it once more.
\end{tcolorbox}
\section[\texorpdfstring{\mw{लाख} (lākh)}{लाख (lākh)}]{\mw{लाख} (lākh) — a hundred thousand}

\label{lesson:MW-C87-lakh}



\begin{quote}
Before the new one: say the Marwadi for ninety, then the Marwadi for a thousand.
\end{quote}

\subsection*{You'll want to know: \mw{लाख}}
\textbf{\mw{लाख}} — \emph{lākh} — \textquotedblleft{}a hundred thousand\textquotedblright{}.

\begin{cousinweb}
\textbf{\mw{लाख}} (\emph{lākh}) is a tadbhava, a word inherited from Sanskrit and reshaped by sound change. It comes from \emph{lakṣa}, \textquotedblleft{}stake, mark; a hundred thousand\textquotedblright{}. CDIAL lists Marwari \emph{lākh} under it (CDIAL 10881).
\end{cousinweb}

\begin{grammarlens}[title={What you've built}]
One more word for this chapter.
\end{grammarlens}

\subsection*{Your turn}
\begin{itemize}
\raggedright
  \item \emph{Say it:} \emph{lākh}
  \item \emph{Say it:} \emph{lākh}, once more
  \item \emph{Say it:} say \emph{lākh}
  \item \emph{Recall:} say the Marwadi for ninety, then the Marwadi for a thousand, then say \emph{lākh} again
\end{itemize}

\begin{tcolorbox}[breakable,colback=teal!4,colframe=teal!35!black,title={Before you move on}]
What does \mw{लाख} mean? (\textquotedblleft{}A hundred thousand\textquotedblright{}.) Say it once more.
\end{tcolorbox}
\section[\texorpdfstring{\mw{दुनिया} (duniyā)}{दुनिया (duniyā)}]{\mw{दुनिया} (duniyā) — the world}

\label{lesson:MW-C87-duniya}



\begin{quote}
Before the new one: say the Marwadi for a thousand, then the Marwadi for a hundred thousand.
\end{quote}

\subsection*{You'll want to know: \mw{दुनिया}}
\textbf{\mw{दुनिया}} — \emph{duniyā} — \textquotedblleft{}the world\textquotedblright{}.

\begin{cousinweb}
\textbf{\mw{दुनिया}} (\emph{duniyā}) is a loan from Arabic \emph{dunyā}, \textquotedblleft{}the world, this life\textquotedblright{}, which reached Marwadi through Persian. Its Arabic root is \emph{d-n-w} (Lane).
\end{cousinweb}

\begin{grammarlens}[title={What you've built}]
One more word for this chapter.
\end{grammarlens}

\subsection*{Your turn}
\begin{itemize}
\raggedright
  \item \emph{Say it:} \emph{duniyā}
  \item \emph{Say it:} \emph{duniyā}, once more
  \item \emph{Say it:} say \emph{duniyā}
  \item \emph{Recall:} say the Marwadi for a thousand, then the Marwadi for a hundred thousand, then say \emph{duniyā} again
\end{itemize}

\begin{tcolorbox}[breakable,colback=teal!4,colframe=teal!35!black,title={Before you move on}]
What does \mw{दुनिया} mean? (\textquotedblleft{}The world\textquotedblright{}.) Say it once more.
\end{tcolorbox}
